\chapter{Systems of Equations}

Solve each system for all variables.

\section{Core Rules}

\begin{rulebox}{Two methods, and when each is faster}
\textbf{Substitution} --- solve one equation for one variable, put it into
the others. Fastest when some equation already has an isolated variable
($x-y=-3$, $L=2K$, $y=x^{2}-4$), and the only method available once the
system is non-linear.

\textbf{Elimination} --- scale equations so that one variable's coefficients
are opposite, then add. Fastest for linear systems with no convenient
isolated variable, and it is the hand version of Gaussian elimination.
\end{rulebox}

\begin{rulebox}{A linear system has exactly three possible outcomes}
\begin{itemize}[leftmargin=1.6em,itemsep=0.3ex]
  \item \textbf{One solution.} The lines cross once.
  \item \textbf{No solution} (inconsistent). Elimination produces a false
        statement such as $0=3$. Geometrically: parallel lines.
  \item \textbf{Infinitely many} (dependent). Elimination produces $0=0$.
        The equations say the same thing; describe the whole solution set
        with a parameter, do not write ``no solution''.
\end{itemize}
There is never any other outcome --- never exactly two. Recognising which
case you are in is itself an examination question.
\end{rulebox}

\begin{rulebox}{Non-linear systems: substitute, then check}
Substitution can raise the degree, and raising the degree manufactures roots
that solve the new equation but not the old one. Squaring both sides of an
equation containing $\sqrt{x}$ is the usual culprit. \textbf{Every candidate
must be substituted back into the original system}, and roots that fail are
discarded --- as are roots outside a domain, such as a negative quantity or
a negative $\sqrt{x}$.
\end{rulebox}

\begin{rulebox}{Lagrangian systems: eliminate $\lambda$, do not solve for it}
First-order conditions arrive in the shape
\[
  \underbrace{U_{x}=\lambda p_{x}}_{(1)},\qquad
  \underbrace{U_{y}=\lambda p_{y}}_{(2)},\qquad
  \underbrace{p_{x}x+p_{y}y=M}_{(3)} .
\]
Divide (1) by (2). The multiplier cancels and what remains is the
\textbf{tangency condition}
\[
  \frac{U_{x}}{U_{y}}=\frac{p_{x}}{p_{y}} ,
\]
one equation in $x$ and $y$ alone. Combine it with the constraint~(3) to
finish. $\lambda$ is recovered at the end from (1) if the question asks for
it, and usually it does not.
\end{rulebox}

\section{Worked Examples}

\begin{worked}{(a) Cobb--Douglas demand: solve $\dfrac{\alpha}{x}=\lambda p_{x}$, $\dfrac{1-\alpha}{y}=\lambda p_{y}$, $p_{x}x+p_{y}y=I$}
\wstep{Step 1. Eliminate $\lambda$ by dividing the first condition by the second.}
\[
  \frac{\alpha/x}{(1-\alpha)/y}=\frac{\lambda p_{x}}{\lambda p_{y}}
  \quad\Longrightarrow\quad
  \frac{\alpha y}{(1-\alpha)x}=\frac{p_{x}}{p_{y}} .
\]
$\lambda$ is gone in one step and never returns. Three equations in three
unknowns have become two in two.
\wtrap{Do not solve the first condition for $\lambda=\alpha/(xp_{x})$ and
substitute. It is not wrong, but it drags an extra symbol through every
later line. Dividing the conditions is the move; it works because $\lambda$
appears to the first power in each.}

\wstep{Step 2. Cross-multiply into a relation between the two expenditures.}
\[
  \alpha\, y p_{y}=(1-\alpha)\,x p_{x}
  \quad\Longrightarrow\quad
  p_{y}y=\frac{1-\alpha}{\alpha}\;p_{x}x .
\]
Read this before continuing: the ratio of spending on the two goods is
$\alpha:(1-\alpha)$, fixed, independent of prices and of income. That is the
signature property of Cobb--Douglas, and it is visible here.

\wstep{Step 3. Substitute into the budget constraint.}
\[
  p_{x}x+\frac{1-\alpha}{\alpha}p_{x}x=I
  \quad\Longrightarrow\quad
  p_{x}x\Bigl(\frac{\alpha+1-\alpha}{\alpha}\Bigr)=I
  \quad\Longrightarrow\quad
  \frac{p_{x}x}{\alpha}=I .
\]
\wtrap{The bracket collapses to $1/\alpha$ because $\alpha+(1-\alpha)=1$.
Grinding through $1+\frac{1-\alpha}{\alpha}$ without noticing this is where
the algebra usually goes wrong.}

\wstep{Step 4. Read off both demands.}
\[
  \boxed{\;x^{*}=\frac{\alpha I}{p_{x}},\qquad
         y^{*}=\frac{(1-\alpha)I}{p_{y}}\;}
\]
the second by symmetry --- or from Step~2, which gives $p_{y}y=(1-\alpha)I$
at once.

\wstep{Step 5. Sanity-check the result.}
Dimensions: income over price is a quantity. \checkmark\ Budget:
$\alpha I+(1-\alpha)I=I$, exactly exhausted. \checkmark\ Signs: $x^{*}>0$
whenever $0<\alpha<1$. \checkmark\ Comparative statics: demand for $x$ falls
in $p_{x}$, rises in $I$, and does not react to $p_{y}$ at all --- the
famous zero cross-price effect.
\end{worked}

\begin{worked}{(b) Two systems that look ordinary and are not: $2x+3y=6,\ 4x+6y=15$ and $x-y=4,\ 3x-3y=12$}
\wstep{Step 1. First system --- eliminate and watch what survives.}
Multiply the first equation by $2$: $4x+6y=12$. Subtract it from the second:
\[
  (4x+6y)-(4x+6y)=15-12
  \quad\Longrightarrow\quad
  0=3 .
\]
A false statement. No pair $(x,y)$ can satisfy both equations, so the system
is \textbf{inconsistent} and the solution set is empty. Geometrically the two
lines have the same slope $-2/3$ but different intercepts: they are parallel
and never meet.

\wstep{Step 2. Second system --- the same move, a different residue.}
The second equation is exactly three times the first, so eliminating gives
\[
  0=0 ,
\]
a statement that is true but empty. The two equations carry one piece of
information between them, and the system is \textbf{dependent}: infinitely
many solutions. Describe them with a parameter,
\[
  \boxed{\;x=t,\quad y=t-4,\quad t\in\R\;}
\]
which is the line $y=x-4$ itself.
\wtrap{$0=0$ means \emph{infinitely many} solutions and $0=3$ means
\emph{none}. Both lines look like ``the variables vanished'', and they are
routinely confused. Read the surviving numbers: a false numeric statement
kills the system, a true one merely says you learned nothing new.}

\wstep{Step 3. In economic terms.}
Two independent markets pin down a unique equilibrium. Two contradictory
requirements --- a price floor above a price ceiling --- give an empty
feasible set. Two restatements of the same requirement leave a whole line of
optima, which is what degeneracy and multiple optima mean in Volume~8.
\end{worked}

\section{Practice Problems}

\begin{problems}
  \item \textbf{[Part A: Linear $2\times 2$ and economic models, 1--10.]}
        $3x+2y=16,\ x-y=-3$
  \item $4x+3y=17,\ 2x-y=1$
  \item $0.5x+0.3y=5.2,\ x+y=12$
  \item Equilibrium: $Q_{d}=80-2P,\ Q_{s}=-10+3P$
  \item IS-LM: $Y=600-200r,\ Y=400+200r$
  \item Break-even: $C=12+3Q,\ R=5Q$, find break-even $Q$.
  \item $2x+3y=6,\ 4x+6y=15$
  \item $x-y=4,\ 3x-3y=12$
  \item $5x-3y=1,\ 2x+y=7$
  \item Capital--labour: $L=2K,\ wL+rK=C$ with $w=3,\,r=4,\,C=50$
  \item \textbf{[Part B: Non-linear $2\times 2$, 11--15.]}
        $y=x^{2}-4,\ y=-x^{2}+4$
  \item $xy=18,\ x-y=3$
  \item $x^{2}+y^{2}=20,\ x=2y$
  \item $y=\mathrm{e}^{x},\ y=\mathrm{e}^{2-x}$
  \item $\sqrt{x}+y=4,\ y=x-2$ (check for extraneous roots)
  \item \textbf{[Part C: Lagrangian-style systems, 16--22.]}
        Cobb--Douglas: $\dfrac{\alpha}{x}=\lambda p_{x},\
        \dfrac{1-\alpha}{y}=\lambda p_{y},\ p_{x}x+p_{y}y=I$, solve
        for $x,y$.
  \item CES capital--labour ratio: $\alpha K^{\rho-1}=\lambda r,\
        (1-\alpha)L^{\rho-1}=\lambda w$, find $K/L$.
  \item $2x=\lambda,\ \lambda=4y,\ 3x+y=14$
  \item $\lambda y=6,\ \lambda x=4,\ xy=24$ ($\lambda>0$)
  \item $4=\lambda K,\ 1=\lambda L,\ K+L=10$
  \item Quasi-linear: $1=\lambda p_{x},\ \dfrac{1}{y}=\lambda p_{y},\
        p_{x}x+p_{y}y=I$, solve for $x,y$.
  \item Simplex: $\alpha x^{\alpha-1}y^{1-\alpha}=\lambda,\
        (1-\alpha)x^{\alpha}y^{-\alpha}=\lambda,\ x+y=1$
  \item \textbf{[Part D: Linear $3\times 3$, 23--30.]}
        $x+y-z=2,\ 2x-y+z=4,\ -x+y+2z=1$
  \item $x+y+z=12,\ 2x+3y-z=13,\ x-y+2z=9$
  \item $x+y+z=2,\ 2x-y-z=1,\ x+2y-z=-3$
  \item $x+y+z=60,\ y=2x,\ z=x+y$
  \item $x+2y+3z=11,\ 4x+y+2z=17,\ 3x+5y+z=16$
  \item $x-y=1,\ y-z=1,\ x+y+z=6$
  \item $x+y+z=3,\ x+2y-z=1,\ -x+y+2z=0$
  \item $2x+y-z=2,\ x-2y+z=0,\ 3x+y+2z=6$
\end{problems}

\section{Solutions}

\begin{solutions}
  \item $x=2,\ y=5$
  \item $x=2,\ y=3$
  \item $x=8,\ y=4$
  \item $P=18,\ Q=44$
  \item $r=\tfrac12,\ Y=500$
  \item $Q=6$
  \item No solution (parallel lines).
  \item Infinitely many solutions (same line).
  \item $x=2,\ y=3$
  \item $K=5,\ L=10$
  \item $(2,0)$ and $(-2,0)$
  \item $(6,3)$ and $(-3,-6)$
  \item $(4,2)$ and $(-4,-2)$
  \item $(1,\mathrm{e})$
  \item $(4,2)$ only; $(9,7)$ is extraneous because
        $\sqrt{9}=3\neq 6-9$.
  \item $x=\dfrac{\alpha I}{p_{x}},\ y=\dfrac{(1-\alpha)I}{p_{y}},\
        \lambda=\dfrac{1}{I}$
  \item $\dfrac{K}{L}=\left(\dfrac{(1-\alpha)r}{\alpha w}\right)^{1/(\rho-1)}$
  \item $x=4,\ y=2,\ \lambda=8$
  \item $\lambda=1,\ x=4,\ y=6$
  \item $\lambda=\tfrac12,\ K=8,\ L=2$
  \item $\lambda=\dfrac{1}{p_{x}},\ y=\dfrac{p_{x}}{p_{y}},\
        x=\dfrac{I}{p_{x}}-1$
  \item $x=\alpha,\ y=1-\alpha$
  \item $(2,1,1)$
  \item $(3,4,5)$
  \item $(1,-1,2)$
  \item $x=10,\ y=20,\ z=30$
  \item $(3,1,2)$
  \item $(3,2,1)$
  \item $(2,0,1)$
  \item $(1,1,1)$
\end{solutions}
